% Einheit 3 von 4 – Rechtwinklige Teildreiecke in Figuren und Vermessung (bank/trigonometrie/e3.jsonl, zusammenbau v0.1)
%% TODO Zweigzeile Teil 1: was hier gelernt wird (Satz in Schülersprache aus dem Einheitstitel) – Einheit 3
\einheitenkopf[e3]{Einheit 3 von 4 · Rechtwinklige Teildreiecke in Figuren und Vermessung}
\zweigzeile{ab Kl. 10 (am Gymnasium ab Kl. 9) · P10 oft}
\setcounter{aufgabe}{22}

%% TODO Titel: Ich-kann-Satz fehlt in der Bank (Platzhalter: Kettenname) – Nr. 23
%% TODO Anweisung: ein Satz über der ersten Teilaufgabe fehlt in der Bank – Nr. 23
\begin{aufgabe}{Teildreiecke und Vermessung}
\swz[3]{Fahre im Trapez das rechtwinklige Teildreieck mit der Höhe $h$ links nach. Beschrifte seine Seiten mit H, G, A – vom Basiswinkel unten links aus.}{\trapez[hoehe=h]{5}{3}{2.5}}
%% TODO Darstellung für schwach fehlt in der Bank (trigonometrie-e3-k1-s1-v1; Abschnitt „Für schwache Schüler“)
\swz{Gleichschenkliges Dreieck: Schenkel $7$ cm, Basiswinkel $63^\circ$. Wie lang ist die Höhe $h$ auf die Grundseite?}{}
%% TODO Darstellung für schwach fehlt in der Bank (trigonometrie-e3-k1-s1-v2; Abschnitt „Für schwache Schüler“)
\swz{Gleichschenkliges Dreieck: Schenkel $9$ cm, Basiswinkel $48^\circ$. Wie lang ist die Höhe $h$ auf die Grundseite?}{}
%% TODO Darstellung für schwach fehlt in der Bank (trigonometrie-e3-k1-s1-v3; Abschnitt „Für schwache Schüler“)
\swz{Gleichschenkliges Dreieck: Schenkel $8$ cm, Basiswinkel $71^\circ$. Wie lang ist die Höhe $h$ auf die Grundseite?}{}
%% TODO Darstellung für schwach fehlt in der Bank (trigonometrie-e3-k1-s1-v4; Abschnitt „Für schwache Schüler“)
\swz{Gleichschenkliges Dreieck: Schenkel $6$ cm, Basiswinkel $57^\circ$. Wie lang ist die Höhe $h$ auf die Grundseite?}{}
%% TODO Darstellung für schwach fehlt in der Bank (trigonometrie-e3-k1-s1-v5; Abschnitt „Für schwache Schüler“)
\swz{Gleichschenkliges Dreieck: Schenkel $4$ cm, Basiswinkel $68^\circ$. Wie lang ist die Höhe $h$ auf die Grundseite?}{}
\swfrage{Was bleibt gleich, was ändert sich?}
%% TODO Darstellung für schwach fehlt in der Bank (trigonometrie-e3-k1-s2-v1; Abschnitt „Für schwache Schüler“)
\swz{Trapez: Höhe $4$ m, Basiswinkel $62^\circ$. Wie lang ist der Schenkel $s$?}{}
%% TODO Darstellung für schwach fehlt in der Bank (trigonometrie-e3-k1-s3-v1; Abschnitt „Für schwache Schüler“)
\swz{Trapez: Schenkel $5$ cm, Basiswinkel $58^\circ$. Wie hoch ist das Trapez?}{}
%% TODO Darstellung für schwach fehlt in der Bank (trigonometrie-e3-k1-s4-v1; Abschnitt „Für schwache Schüler“)
\swz{Rechtwinkliges Trapez: die senkrechten Seiten sind $9$ cm und $4$ cm lang und stehen $8$ cm auseinander. Wie groß ist der Winkel $\alpha$ zwischen der schrägen Seite und der Waagerechten?}{}
%% TODO Darstellung für schwach fehlt in der Bank (trigonometrie-e3-k1-s5-v1; Abschnitt „Für schwache Schüler“)
\swz{Parallelogramm $ABCD$: $BC = 9$ cm. $F$ liegt auf der Verlängerung von $AB$ über $B$ hinaus, $CF$ steht senkrecht darauf; der Winkel $CBF$ beträgt $58^\circ$. Wie lang ist die Höhe $h = CF$?}{}
%% TODO Darstellung für schwach fehlt in der Bank (trigonometrie-e3-k1-s6-v1; Abschnitt „Für schwache Schüler“)
\swz{Parallelogramm $ABCD$ mit der Höhe $h = 6$ cm von $C$ auf die Verlängerung von $AB$. Die Diagonale $AC$ bildet mit $AB$ einen Winkel von $27^\circ$. Wie lang ist $AC$?}{}
%% TODO Darstellung für schwach fehlt in der Bank (trigonometrie-e3-k1-s7-v1; Abschnitt „Für schwache Schüler“)
\swz{Drachen: eine obere Seite ist $6$ cm lang und bildet mit der Längsdiagonale einen Winkel von $32^\circ$. Wie lang ist die halbe Querdiagonale?}{}
%% TODO Darstellung für schwach fehlt in der Bank (trigonometrie-e3-k1-s8-v1; Abschnitt „Für schwache Schüler“)
\swz{Aus $45$ m waagerechtem Abstand sieht man die Spitze eines Baums unter einem Höhenwinkel von $31^\circ$. Wie hoch liegt sie über dem Messgerät?}{}
%% TODO Darstellung für schwach fehlt in der Bank (trigonometrie-e3-k1-s9-v1; Abschnitt „Für schwache Schüler“)
\swz{Ein Messgerät steht in $1{,}4$ m Höhe, $58$ m waagerecht entfernt. Der Höhenwinkel zur Spitze eines Kirchturms beträgt $27^\circ$. Wie hoch ist es?}{}
%% TODO Darstellung für schwach fehlt in der Bank (trigonometrie-e3-k1-s10-v1; Abschnitt „Für schwache Schüler“)
\swz{Aus $48$ m Abstand peilt man in $1{,}7$ m Augenhöhe die Dachkante eines Turms unter $24^\circ$ an. Auf dem Dach steht eine $6$ m hohe Antenne. Wie hoch liegt die Antennenspitze über dem Boden?}{}
%% TODO Darstellung für schwach fehlt in der Bank (trigonometrie-e3-k1-s11-v1; Abschnitt „Für schwache Schüler“)
\swz{Rechter Winkel bei $D$; $B$ liegt zwischen $D$ und $C$. Von $A$ aus: $\angle DAC = 71^\circ$, $\angle BAC = 23^\circ$, $AD = 32$ m. Wie lang ist $BD$?}{}
%% TODO Darstellung für schwach fehlt in der Bank (trigonometrie-e3-k1-s12-v1; Abschnitt „Für schwache Schüler“)
\swz{Dreieck $ABC$ mit der Höhe $h$ von $C$ auf $AB$: $AC = 8$ cm, $\alpha = 42^\circ$, $\beta = 61^\circ$. Berechne erst $h$, dann $BC$.}{}
%% TODO Darstellung für schwach fehlt in der Bank (trigonometrie-e3-k1-s13-v1; Abschnitt „Für schwache Schüler“)
\swz[3]{Gleichschenkliges Trapez: bekannt sind die beiden parallelen Seiten und die Höhe, gesucht ist der Umfang. Beschreibe den Rechenweg ohne zu rechnen.}{}
%% TODO Darstellung für schwach fehlt in der Bank (trigonometrie-e3-k1-s14-v1; Abschnitt „Für schwache Schüler“)
\swz{Hypotenuse $9$ cm, $\alpha = 33^\circ$. Berechne Gegenkathete und Ankathete und prüfe mit dem Satz des Pythagoras.}{}
%% TODO Darstellung für schwach fehlt in der Bank (trigonometrie-e3-k1-s15-v1; Abschnitt „Für schwache Schüler“)
\swz{Dreieck $ABC$: $c = 9$ cm, $b = 6$ cm, $\alpha = 47^\circ$. Berechne die Höhe $h$ auf $c$ und dann den Flächeninhalt.}{}
%% TODO Darstellung für schwach fehlt in der Bank (trigonometrie-e3-k1-s16-v1; Abschnitt „Für schwache Schüler“)
\swz{Quadratische Pyramide: Grundkante $6$ m, die Seitenflächen sind um $57^\circ$ gegen die Grundfläche geneigt. Wie lang ist die Seitenhöhe $h_s$?}{}
\swz[3]{Ein Messgerät steht $1{,}4$ m über dem Boden. Die Sichtlinie zur Hügelkuppe ist $96{,}5$ m lang und steigt unter $27^\circ$ an. Wie hoch ist der Hügel? Auf der Kuppe steht ein $9$ m hoher Aussichtsplattform – wie hoch über dem Boden liegt seine Oberkante? \hfill (P10 2018 OS)}{}
\end{aufgabe}

%% TODO Titel: Ich-kann-Satz fehlt in der Bank (Platzhalter: Kettenname) – Nr. 24
%% TODO Anweisung: ein Satz über der ersten Teilaufgabe fehlt in der Bank – Nr. 24
\begin{aufgabe}{Strecke aus Teilstrecken nach der Trigonometrie}
%% TODO Darstellung für schwach fehlt in der Bank (trigonometrie-e3-k2-s1-v1; Abschnitt „Für schwache Schüler“)
\swz{Dreieck $ABC$ mit der Höhe $CD = 4$ cm auf $AB$: $AD = 6$ cm, $\beta = 53^\circ$. Wie lang ist $AB$?}{}
\end{aufgabe}

%% TODO Titel: Ich-kann-Satz fehlt in der Bank (Platzhalter: Kettenname) – Nr. 25
%% TODO Anweisung: ein Satz über der ersten Teilaufgabe fehlt in der Bank – Nr. 25
\begin{aufgabe}{Teildreiecke und Vermessung – Fehler finden, Begründen, Anwendung}
\swz[3]{Abstand zum Turm $44$ m, Höhenwinkel $31^\circ$, das Messgerät steht $1{,}7$ m hoch. Wo ist der Fehler? \rechnung{h &= 44 \cdot \mathrm{tan}\,31^\circ \\ h &\approx 26{,}4 \text{ m}}}{}
\swfrage{Warum rechnet man im gleichschenkligen Dreieck mit der halben Grundseite, wenn man den Basiswinkel benutzt?}
\swz[3]{Ein Baum soll gefällt werden, das Haus steht $18$ m vom Stamm entfernt. Aus $20$ m Abstand sieht man die Baumspitze in $1{,}7$ m Augenhöhe unter $41^\circ$. Kann der Baum beim Fallen das Haus treffen?}{}
\end{aufgabe}

\uebersichtskasten{Teildreieck: Suche das rechtwinklige Dreieck in der Figur, fahre es nach und benenne die Seiten vom gegebenen Winkel aus – die Höhe, ein Abstand oder eine Hilfslinie ist meist eine Kathete, die schräge Seite die Hypotenuse. \\ Trapez mit Höhe 5 m und Basiswinkel 50°: Schenkel = 5 : sin 50° ≈ 6,5 m \\ Vermessung: Abstand und Höhenwinkel geben die Höhe über dem Gerät; Gerätehöhe, Sockel oder Aufbau werden danach addiert. \\ Abstand 40 m, Höhenwinkel 35°, Augenhöhe 1,6 m: Turmhöhe = 40 · tan 35° + 1,6 ≈ 28,0 + 1,6 = 29,6 m \\ Zwei Dreiecke: Das Ergebnis des ersten Dreiecks ist die gegebene Seite des zweiten; ein Teilwinkel ist die Differenz zweier Winkel an derselben Ecke.}
